\chapter{Circuitos RC y RL}\label{ch:g12:rc-rl-circuits}

Pulsa el disparador de una cámara hasta la mitad: un zumbido agudo sube durante dos
segundos, una luz se pone verde --- y solo entonces disparará el flash, soltando en un
milisegundo lo que ha tardado dos segundos en reunir. Este capítulo presenta los dos
componentes que le dan a un circuito memoria y reloj --- el \omterm{def:g12:rc-rl-circuits:capacitor}{condensador} y la \omterm{def:g12:rc-rl-circuits:inductor}{bobina}
--- y las primeras ecuaciones diferenciales de la física, comprobables con las
derivadas de este curso.

\section{El condensador}

\begin{definition}[Condensador y capacidad]\label{def:g12:rc-rl-circuits:capacitor}
Un \emph{condensador}\index{condensador} es un par de placas metálicas enfrentadas
separadas por un aislante delgado. Conectadas a un generador, las placas adquieren
cargas opuestas $+q$ y $-q$ --- un \omterm{def:g11:electric-gravitational-fields:efield}{campo eléctrico} llena el hueco
(\cref{ch:g11:electric-gravitational-fields}) --- y la tensión $u$ entre las placas es
proporcional a la carga almacenada:
\[ q = C\,u . \]
La constante $C$ es la \emph{capacidad}\index{capacidad}, en
\emph{faradios}\index{faradio} (\unit{F}): $\qty{1}{F} = \qty{1}{C/V}$.
\end{definition}

\begin{example}[Órdenes de magnitud]\label{ex:g12:rc-rl-circuits:orders}
Un \omterm{def:g12:rc-rl-circuits:capacitor}{faradio} es enorme --- un \omterm{def:g11:fundamental-interactions:charge}{culombio} aparcado por \omterm{def:g11:circuits-and-power:voltage}{voltio}. Los cerámicos de los
circuitos de radio: de \unit{pF} a \unit{nF}; los electrolíticos de las fuentes de
alimentación y de los flashes: de \unit{\micro F} a \unit{mF}; los supercondensadores:
miles de \omterm{def:g12:rc-rl-circuits:capacitor}{faradios}.
\end{example}

\begin{proposition}[Corriente hacia un condensador]\label{prop:g12:rc-rl-circuits:cdudt}
La corriente que llega a la placa de un \omterm{def:g12:rc-rl-circuits:capacitor}{condensador} es la rapidez con la que crece su
carga:
\[
i = \frac{dq}{dt} = C\,\frac{du}{dt} ,
\]
donde las minúsculas señalan magnitudes variables en el tiempo ($U$, $I$ se reservan
para las constantes, \cref{ch:g11:circuits-and-power}).
\end{proposition}

\begin{proof}
La corriente es la carga entregada por segundo
(\cref{def:g11:circuits-and-power:current}): sobre un intervalo cada vez más pequeño,
$i$ es la derivada de $q$; y $q = Cu$ con $C$ constante da $C\,du/dt$.
\end{proof}

\begin{remark}[Dos consecuencias]\label{rem:g12:rc-rl-circuits:continuity}
Si $u$ es constante, $i = 0$: un \omterm{def:g12:rc-rl-circuits:capacitor}{condensador} cargado no deja pasar ninguna corriente
continua. Y $u$ no puede dar nunca un salto --- un escalón exigiría una corriente
infinita: la tensión del \omterm{def:g12:rc-rl-circuits:capacitor}{condensador} es la memoria del circuito.
\end{remark}

\section{Cargar un condensador: el circuito RC}

Una sola malla --- generador $\mathcal{E}$
(\cref{def:g11:circuits-and-power:emf}), interruptor, resistor y \omterm{def:g12:rc-rl-circuits:capacitor}{condensador}
descargado: se cierra $K$ en $t = 0$.

\begin{omfigure}
\begin{circuitikz}[font=\small]
  \draw (0,0) to[battery1, l=$\mathcal{E}$] (0,2.4) to[nos, l=$K$] (2.1,2.4)
    to[R=$R$, i=$i$] (4.2,2.4) -- (5.4,2.4) to[C=$C$, v=$u$] (5.4,0) -- (0,0);
\end{circuitikz}

{\small El \omterm{thm:g12:rc-rl-circuits:charging}{circuito RC} de carga: al cerrar $K$, el generador empuja carga a través de
$R$ hasta las placas de $C$.}
\end{omfigure}

\begin{theorem}[Ley de carga del circuito RC]\label{thm:g12:rc-rl-circuits:charging}
\index{circuito RC}
Tras cerrarse el interruptor, la tensión del \omterm{def:g12:rc-rl-circuits:capacitor}{condensador} obedece la ecuación siguiente
y, a partir de $u(0) = 0$, resulta:
\[
RC\,\frac{du}{dt} + u = \mathcal{E},
\qquad
u(t) = \mathcal{E}\left(1 - e^{-t/\tau}\right),
\qquad \tau = RC .
\]
\end{theorem}

\begin{proof}
Las tensiones a lo largo de la malla se suman
(\cref{prop:g11:circuits-and-power:laws}): $\mathcal{E} = u_R + u = Ri + u$ en todo
instante; sustitúyase $i = C\,du/dt$ (\cref{prop:g12:rc-rl-circuits:cdudt}). Ahora
\emph{verifiquemos} la candidata: $du/dt = (\mathcal{E}/\tau)e^{-t/\tau}$, luego
$RC\,(\mathcal{E}/RC)\,e^{-t/\tau} + \mathcal{E}(1 - e^{-t/\tau}) = \mathcal{E}$,
cierto para toda $t$; y $u(0) = \mathcal{E}(1-1) = 0$. Que ninguna \emph{otra} curva
encaja a la vez con la ecuación y con el arranque se demuestra en el curso de
matemáticas de este año, al tratar las ecuaciones diferenciales.
\end{proof}

\begin{definition}[Constante de tiempo]\label{def:g12:rc-rl-circuits:tau}
La \emph{constante de tiempo}\index{constante de tiempo} de un \omterm{thm:g12:rc-rl-circuits:charging}{circuito RC} es
$\tau = RC$. \omterm{def:g10:orders-of-magnitude:unit}{Unidades}: $\unit{\ohm} \times \unit{F} = (\unit{V/A})(\unit{C/V}) =
\unit{C/A} = \unit{s}$ --- el tempo propio del circuito.
\end{definition}

\begin{method}[Leer $\tau$ en la curva]\label{met:g12:rc-rl-circuits:reading}
Tres lecturas equivalentes, directamente sobre la gráfica:
\begin{enumerate}
  \item \emph{tangente en el origen}: $u'(0) = \mathcal{E}/\tau$, así que la tangente
        inicial corta la asíntota en $t = \tau$;
  \item \emph{63\,\%}: $u(\tau) = \mathcal{E}(1 - e^{-1}) \approx 0.63\,\mathcal{E}$;
  \item \emph{95\,\%}: $u(3\tau) \approx 0.95\,\mathcal{E}$ --- cargado, a efectos
        prácticos ($5\tau$: $99\,\%$).
\end{enumerate}
\end{method}

\begin{proposition}[Descarga]\label{prop:g12:rc-rl-circuits:discharge}
Un \omterm{def:g12:rc-rl-circuits:capacitor}{condensador} cargado a $U_0$ y cerrado únicamente sobre un resistor $R$ obedece
$RC\,du/dt + u = 0$ y sigue
\[
u(t) = U_0\,e^{-t/\tau}, \qquad \tau = RC
\]
--- queda el $37\,\%$ en $\tau$ y el $5\,\%$ en $3\tau$.
\end{proposition}

\begin{proof}
La misma ley de mallas, sin generador; y de nuevo por sustitución:
$RC(-U_0/\tau)e^{-t/\tau} + U_0 e^{-t/\tau} = 0$, con $u(0) = U_0$.
\end{proof}

\begin{omfigure}
\begin{tikzpicture}
\begin{axis}[omaxis, width=9.2cm, height=5.4cm, xlabel={$t$}, ylabel={$u$},
    xmin=0, xmax=5, ymin=0, ymax=7.4, xtick={0,1,2,3,4,5},
    xticklabels={$0$,$\tau$,$2\tau$,$3\tau$,$4\tau$,$5\tau$},
    ytick={0,2.21,3.79,6.0},
    yticklabels={$0$,$0.37\,\mathcal{E}$,$0.63\,\mathcal{E}$,$\mathcal{E}$}]
  \addplot[omExo, dashed, domain=0:5, samples=2] {6};
  \addplot[omExo, dashed, domain=0:1.18, samples=2] {6*x};
  \addplot[omExo, dashed, domain=0:1.1, samples=2] {6-6*x};
  \addplot[omDef, thick, domain=0:5, samples=80] {6*(1-exp(-x))};
  \addplot[omThm, thick, domain=0:5, samples=80] {6*exp(-x)};
  \draw[dashed, omExo] (1,0) -- (1,3.79) -- (0,3.79);
  \draw[dashed, omExo] (3,0) -- (3,5.70);
\end{axis}
\end{tikzpicture}

{\small Carga (azul, desde $u = 0$ hacia $\mathcal{E}$) y descarga (rojo, desde
$U_0 = \mathcal{E}$). Las dos tangentes iniciales cortan su asíntota en $t = \tau$; el
$63\,\%$ está hecho en $\tau$ y el $95\,\%$ en $3\tau$.}
\end{omfigure}

\begin{example}[Los números de un flash]\label{ex:g12:rc-rl-circuits:flashrc}
Un flash fotográfico carga $C = \qty{800}{\micro F}$ a través de
$R = \qty{1.0}{k\ohm}$: $\tau = 10^3 \times 8.0\times10^{-4} = \qty{0.80}{s}$, así que
la luz de ``listo'' --- cableada para saltar al $95\,\%$ --- espera
$3\tau \approx \qty{2.4}{s}$: el zumbido del enganche. Al dispararse a través de los
meros \qty{0.50}{\ohm} del tubo de descarga, el \emph{mismo} \omterm{def:g12:rc-rl-circuits:capacitor}{condensador} se vacía con
$\tau' = \qty{0.40}{ms}$: dos mil veces más rápido para salir que para entrar.
\end{example}

\section{La bobina: autoinducción y circuito RL}

\begin{definition}[Bobina y autoinducción]\label{def:g12:rc-rl-circuits:inductor}
Una \emph{bobina}\index{bobina} es un arrollamiento de hilo: su corriente enhebra sus
propias espiras con un \omterm{def:g11:magnetic-fields:bfield}{campo magnético} (\cref{ch:g11:magnetic-fields}); y cada vez que
este \emph{cambia}, aparece una tensión entre los bornes de la bobina:
\[
u = L\,\frac{di}{dt} .
\]
La constante $L$ es la \emph{autoinducción}\index{autoinducción}, en
\emph{henrios}\index{henrio} (\unit{H}): $\qty{1}{H} = \qty{1}{V.s/A}$.
\end{definition}
\admitted

\begin{remark}[La bobina se resiste al cambio]\label{rem:g12:rc-rl-circuits:oppose}
Por qué una corriente variable genera una tensión --- la inducción --- se deduce en
los volúmenes universitarios; aquí tomamos la ley y su carácter: la \omterm{def:g12:rc-rl-circuits:inductor}{bobina} \emph{se
opone a los cambios} de su corriente. Gemelas en espejo: la $u$ de un \omterm{def:g12:rc-rl-circuits:capacitor}{condensador} no
puede saltar; la $i$ de una \omterm{def:g12:rc-rl-circuits:inductor}{bobina}, tampoco.
\end{remark}

\begin{omfigure}
\begin{circuitikz}[font=\small]
  \draw (0,0) to[battery1, l=$\mathcal{E}$] (0,2.4) to[nos, l=$K$] (2.1,2.4)
    to[R=$R$] (4.2,2.4) -- (5.4,2.4) to[L=$L$, i=$i$, v=$u$] (5.4,0) -- (0,0);
\end{circuitikz}

{\small El \omterm{prop:g12:rc-rl-circuits:rl}{circuito RL}: tras cerrarse $K$, la \omterm{def:g12:rc-rl-circuits:inductor}{bobina} deja subir la corriente solo a su
propio ritmo $\tau = L/R$.}
\end{omfigure}

\begin{proposition}[Subida de la corriente en un circuito RL]\label{prop:g12:rc-rl-circuits:rl}
\index{circuito RL}
Al cerrar el interruptor sobre un generador $\mathcal{E}$, un resistor $R$ y una
\omterm{def:g12:rc-rl-circuits:inductor}{bobina} $L$ en serie, resulta $\mathcal{E} = Ri + L\,di/dt$, es decir,
\[
\frac{L}{R}\,\frac{di}{dt} + i = \frac{\mathcal{E}}{R},
\qquad\text{resuelta por}\quad
i(t) = \frac{\mathcal{E}}{R}\left(1 - e^{-t/\tau}\right),
\quad \tau = \frac{L}{R} .
\]
\end{proposition}

\begin{proof}
De nuevo la ley de mallas; y es la \emph{misma ecuación} que en el
\cref{thm:g12:rc-rl-circuits:charging} con $u \to i$, $\mathcal{E} \to \mathcal{E}/R$
y $RC \to L/R$ --- la misma sustitución verifica la solución, y todas las lecturas del
\cref{met:g12:rc-rl-circuits:reading} se trasladan sin cambios.
\end{proof}

\begin{example}[Un electroimán se despierta]\label{ex:g12:rc-rl-circuits:rl}
La \omterm{def:g12:rc-rl-circuits:inductor}{bobina} de un relé, $L = \qty{0.30}{H}$ y $R = \qty{6.0}{\ohm}$, sobre una
alimentación de \qty{12}{V}: corriente final $\mathcal{E}/R = \qty{2.0}{A}$, al ritmo
$\tau = 0.30/6.0 = \qty{50}{ms}$ --- despierto en
$3\tau \approx \qty{0.15}{s}$.
\end{example}

\begin{remark}[La chispa al abrir]\label{rem:g12:rc-rl-circuits:spark}
Abre un interruptor sobre una \omterm{def:g12:rc-rl-circuits:inductor}{bobina} en servicio y la $i$ tiene que desplomarse a cero
en una fracción de milisegundo: $di/dt$ es enorme y $u = L\,di/dt$ alcanza cientos o
miles de \omterm{def:g11:circuits-and-power:voltage}{voltios} --- una chispa salta el contacto. La \omterm{def:g12:rc-rl-circuits:inductor}{bobina} de encendido de un coche
lo hace a propósito, miles de veces por minuto.
\end{remark}

\section{Energía almacenada}

\begin{proposition}[Energía en un condensador, energía en una bobina]\label{prop:g12:rc-rl-circuits:energy}
\index{energía!de un condensador}\index{energía!de una bobina}
Un \omterm{def:g12:rc-rl-circuits:capacitor}{condensador} $C$ cargado a la tensión $u$ y una \omterm{def:g12:rc-rl-circuits:inductor}{bobina} $L$ recorrida por una
corriente $i$ almacenan
\[
E_C = \tfrac12\,C u^2 = \frac{q^2}{2C},
\qquad
E_L = \tfrac12\,L i^2 .
\]
\end{proposition}

\begin{proof}
Depositar una carga pequeña $dq$ sobre unas placas que ya están a la tensión $u$
cuesta $u\,dq$ (\cref{def:g11:circuits-and-power:voltage}); apilando los costes, el
total es el área bajo la recta $u = q/C$ --- el triángulo de abajo ---, luego
$E_C = \tfrac12 qu = \tfrac12 Cu^2$. El mismo triángulo para la \omterm{def:g12:rc-rl-circuits:inductor}{bobina} (subir $i$ en
$di$ cuesta $Li\,di$) da $E_L = \tfrac12 Li^2$.
\end{proof}

\begin{omfigure}
\begin{tikzpicture}[font=\small, x=1.35cm, y=1.05cm]
  \fill[omDef!25] (0,0) -- (4,0) -- (4,2.6) -- cycle;
  \fill[omDef!60] (2.4,0) rectangle (2.7,1.658);
  \draw[omDef, thick] (0,0) -- (4,2.6);
  \draw[-Stealth] (0,0) -- (4.8,0) node[right] {$q$};
  \draw[-Stealth] (0,0) -- (0,3.1) node[above] {$u = q/C$};
  \draw[dashed] (4,0) -- (4,2.6) -- (0,2.6);
  \node[below] at (4,0) {$q_0$}; \node[left] at (0,2.6) {$u_0$};
  \node[below, font=\footnotesize] at (2.55,0) {$dq$};
  \node[omDef!50!black, font=\footnotesize] at (2.9,0.85)
    {$E_C$};
\end{tikzpicture}

{\small Cada rebanada $dq$ cuesta $u\,dq$; la carga entera cuesta el área del
triángulo, $E_C = \tfrac12 q_0 u_0$ --- los primeros \omterm{def:g11:fundamental-interactions:charge}{culombios} embarcan a tensión
baja.}
\end{omfigure}

\begin{example}[Órdenes de magnitud de la energía]\label{ex:g12:rc-rl-circuits:energies}
El \omterm{def:g12:rc-rl-circuits:capacitor}{condensador} del flash del \cref{ex:g12:rc-rl-circuits:flashrc}:
$\tfrac12 \times 8.0\times10^{-4} \times 300^2 = \qty{36}{J}$. Un supercondensador de
\qty{3000}{F} a \qty{2.7}{V}: \qty{1.1e4}{J} --- lo que da una pila AA. Una \omterm{def:g12:rc-rl-circuits:inductor}{bobina} de
\qty{10}{mH} a \qty{10}{A}: \qty{0.50}{J}. Menos de lo que guardan las pilas --- pero
liberable en un milisegundo.
\end{example}

\begin{remark}[Para qué sirven los dos guardianes del tiempo]\label{rem:g12:rc-rl-circuits:applications}
\index{alisado}\index{circuito temporizador}
Entrar despacio y salir deprisa es lo que mueve el flash de la cámara y el
desfibrilador. El retardo $\tau = RC$ hace funcionar los \omterm{rem:g12:rc-rl-circuits:applications}{circuitos temporizadores}:
limpiaparabrisas intermitentes, balizas que parpadean, luces de escalera. Un
\omterm{def:g12:rc-rl-circuits:capacitor}{condensador} colocado en paralelo con una alimentación \emph{alisa} la tensión ($u$ no
puede saltar); una \omterm{def:g12:rc-rl-circuits:inductor}{bobina} en serie alisa la corriente ($i$ no puede saltar). Y \omterm{def:g12:rc-rl-circuits:inductor}{bobina}
más \omterm{def:g12:rc-rl-circuits:capacitor}{condensador} juntos forman un \omterm{def:g12:oscillators-and-time:oscillator}{oscilador} (\cref{ch:g12:rlc-oscillations}).
\end{remark}

\section{Ejercicios}

\begin{exercise}[$\star$]\label{exo:g12:rc-rl-circuits:1}
Un \omterm{def:g12:rc-rl-circuits:capacitor}{condensador} de \qty{470}{\micro F} se carga a \qty{12}{V}: ¿qué carga alberga? ¿Qué
tensión pondría \qty{2.0}{mC} en \qty{100}{\micro F}?
\end{exercise}

\begin{exercise}[$\star$]\label{exo:g12:rc-rl-circuits:2}
Reparte \qty{47}{pF}, \qty{2200}{\micro F} y \qty{10}{F} entre un alisador de fuente
de alimentación, un sintonizador de radio y un supercondensador; calcula la \omterm{def:g10:energy-conservation:energy}{energía}
del último a \qty{2.7}{V}.
\end{exercise}

\begin{exercise}[$\star$]\label{exo:g12:rc-rl-circuits:3}
La tensión entre los bornes de un \omterm{def:g12:rc-rl-circuits:capacitor}{condensador} de \qty{220}{\micro F} sube de forma
uniforme de $0$ a \qty{5.0}{V} en \qty{10}{ms}: ¿qué corriente entra? ¿Y qué corriente
una vez que la tensión se queda quieta, y por qué?
\end{exercise}

\begin{exercise}[$\star$]\label{exo:g12:rc-rl-circuits:4}
Calcula $\tau$ para $R = \qty{10}{k\ohm}$ y $C = \qty{100}{\micro F}$, y después para
$L = \qty{0.10}{H}$ y $R = \qty{50}{\ohm}$. Comprueba, \omterm{def:g10:orders-of-magnitude:unit}{unidad} a \omterm{def:g10:orders-of-magnitude:unit}{unidad}, que ambas
combinaciones son segundos.
\end{exercise}

\begin{exercise}[$\star$]\label{exo:g12:rc-rl-circuits:5}
Un \omterm{def:g12:rc-rl-circuits:capacitor}{condensador} se carga hacia \qty{6.0}{V} y pasa por \qty{3.8}{V} en
$t = \qty{2.0}{s}$. Deduce la \omterm{def:g12:rc-rl-circuits:tau}{constante de tiempo}; ¿cuándo está la carga hecha al
$95\,\%$?
\end{exercise}

\begin{exercise}[$\star\star$]\label{exo:g12:rc-rl-circuits:6}
Un \omterm{thm:g12:rc-rl-circuits:charging}{circuito RC}: $\mathcal{E} = \qty{9.0}{V}$, $R = \qty{4.7}{k\ohm}$,
$C = \qty{220}{\micro F}$. Calcula $\tau$, $u(\tau)$, $u(3\tau)$ y la corriente en el
instante en que se cierra el interruptor.
\end{exercise}

\begin{exercise}[$\star\star$]\label{exo:g12:rc-rl-circuits:7}
Verifica por sustitución que $u(t) = \mathcal{E}(1 - e^{-t/RC})$ satisface
$RC\,du/dt + u = \mathcal{E}$; ¿qué codifican físicamente su valor en $t = 0$ y su
límite para $t$ grande?
\end{exercise}

\begin{exercise}[$\star\star$]\label{exo:g12:rc-rl-circuits:8}
Un \omterm{def:g12:rc-rl-circuits:capacitor}{condensador} de \qty{100}{\micro F} cargado a \qty{12}{V} se descarga a través de
\qty{47}{k\ohm}: calcula $\tau$, $u(\tau)$ y el tiempo que tarda $u$ en bajar de
\qty{0.60}{V}.
\end{exercise}

\begin{exercise}[$\star\star$]\label{exo:g12:rc-rl-circuits:9}
La corriente en una \omterm{def:g12:rc-rl-circuits:inductor}{bobina} de \qty{0.50}{H} sube de forma uniforme de $0$ a
\qty{2.0}{A} en \qty{40}{ms}: ¿qué tensión aparece entre sus bornes? Con
$u = L\,di/dt$, ¿por qué salta la chispa al \emph{abrir} un interruptor sobre una
\omterm{def:g12:rc-rl-circuits:inductor}{bobina} y no al cerrarlo?
\end{exercise}

\begin{exercise}[$\star\star$]\label{exo:g12:rc-rl-circuits:10}
El relé del \cref{ex:g12:rc-rl-circuits:rl}: verifica por sustitución que
$i(t) = (\mathcal{E}/R)(1 - e^{-Rt/L})$ satisface $L\,di/dt + Ri = \mathcal{E}$ y
calcula después $i(\tau)$ e $i(3\tau)$.
\end{exercise}

\begin{exercise}[$\star\star$]\label{exo:g12:rc-rl-circuits:11}
Demuestra que $u'(0) = \mathcal{E}/\tau$ y deduce que la tangente en el origen corta
la asíntota en $t = \tau$. Sobre una curva registrada que da $\tau = \qty{0.50}{s}$,
con $R = \qty{2.0}{k\ohm}$: halla $C$.
\end{exercise}

\begin{exercise}[$\star\star\star$]\label{exo:g12:rc-rl-circuits:12}
Ordena por \omterm{def:g10:energy-conservation:energy}{energía} almacenada, calculando cada una: un supercondensador de
\qty{3000}{F} a \qty{2.7}{V}; un \omterm{def:g12:rc-rl-circuits:capacitor}{condensador} de flash, \qty{800}{\micro F} a
\qty{300}{V}; una \omterm{def:g12:rc-rl-circuits:inductor}{bobina} de \qty{10}{mH} a \qty{10}{A}; una pila AA
($E \approx \qty{1.3e4}{J}$). ¿Qué ofrecen los tres primeros que la pila no puede dar?
\end{exercise}

\begin{exercise}[$\star\star\star$]\label{exo:g12:rc-rl-circuits:13}
Un temporizador de escalera mantiene la luz encendida mientras
$u < 0.63\,\mathcal{E}$ entre los bornes de su \omterm{def:g12:rc-rl-circuits:capacitor}{condensador} en carga, y se apaga en
$t = \tau$. Diseña la luz para que dure \qty{5.0}{s} usando $C = \qty{100}{\micro F}$:
¿qué $R$? Para un umbral $0.50\,\mathcal{E}$, demuestra que el retardo vale
$\tau \ln 2$ y calcúlalo.
\end{exercise}

\begin{exercise}[$\star\star\star$]\label{exo:g12:rc-rl-circuits:14}
Un \omterm{def:g12:rc-rl-circuits:capacitor}{condensador} de \qty{2200}{\micro F} alisa una alimentación: entre dos recargas, él
solo alimenta una carga que consume \qty{0.50}{A} durante \qty{10}{ms}. A partir de
$i = C\,du/dt$, calcula la caída de tensión y después la \omterm{def:g12:rc-rl-circuits:capacitor}{capacidad} que la mantiene por
debajo de \qty{0.50}{V}.
\end{exercise}

\begin{exercise}[$\star\star\star$]\label{exo:g12:rc-rl-circuits:15}
Una \omterm{def:g12:rc-rl-circuits:inductor}{bobina} de encendido, $L = \qty{0.80}{H}$, lleva \qty{2.0}{A} cuando el ruptor la
corta en aproximadamente \qty{1.0}{ms}. Estima la tensión de la \omterm{def:g12:rc-rl-circuits:inductor}{bobina} durante el
corte y la \omterm{def:g10:energy-conservation:energy}{energía} de la chispa; ¿por qué una interrupción lenta no da chispa?
\end{exercise}

\section{Problema: El flash de la cámara}

\begin{problem}[{Problema de fin de semana --- el flash de la cámara: dos segundos de
zumbido, un milisegundo de gloria y el primo de pared que reinicia corazones --- una
sola exponencial los gobierna a todos}]
\label{pb:g12:rc-rl-circuits:1}
Un flash compacto almacena \omterm{def:g10:energy-conservation:energy}{energía} en un \omterm{def:g12:rc-rl-circuits:capacitor}{condensador} $C = \qty{800}{\micro F}$,
cargado a través de $R = \qty{1.0}{k\ohm}$ desde un \omterm{def:g10:energy-conservation:transfer}{convertidor} modelizado como un
generador ideal $\mathcal{E} = \qty{300}{V}$; la luz de ``listo'' salta al $95\,\%$ de
la tensión final. Al dispararse, el \omterm{def:g12:rc-rl-circuits:capacitor}{condensador} solo ve el tubo de descarga,
$R_f = \qty{0.50}{\ohm}$. Datos: una pila AA entrega \qty{1.5}{V} y \qty{2.4}{A.h};
$e = \qty{1.60e-19}{C}$.

\medskip
\noindent\textbf{Parte I --- El depósito.}
\begin{enumerate}
  \item Calcula la carga del \omterm{def:g12:rc-rl-circuits:capacitor}{condensador} a \qty{300}{V}.
  \item ¿A cuántas \omterm{def:g11:fundamental-interactions:elementary}{cargas elementales} equivale?
  \item Calcula la \omterm{def:g10:energy-conservation:energy}{energía} almacenada $E_C$.
  \item La pila AA alberga $q = It \approx \qty{8.6e3}{C}$ y, por tanto, unos
        \qty{1.3e4}{J}. ¿Cuántas veces la \omterm{def:g10:energy-conservation:energy}{energía} del \omterm{def:g12:rc-rl-circuits:capacitor}{condensador} es eso?
  \item Entonces, ¿para qué molestarse con el \omterm{def:g12:rc-rl-circuits:capacitor}{condensador}? Una frase, usando la
        palabra ``\omterm{def:g11:work-of-force:power}{potencia}''.
\end{enumerate}

\medskip
\noindent\textbf{Parte II --- La espera del verde.}
\begin{enumerate}[resume]
  \item Escribe la ley de mallas y conviértela en la ecuación diferencial de $u(t)$.
  \item Verifica por sustitución que $u(t) = \mathcal{E}(1 - e^{-t/RC})$ la resuelve
        partiendo de $u(0) = 0$.
  \item Calcula la \omterm{def:g12:rc-rl-circuits:tau}{constante de tiempo} $\tau$.
  \item ¿Al cabo de cuánto tiempo se pone verde la luz de listo, y a qué tensión?
        Reconoce el zumbido de dos segundos.
  \item Calcula la corriente inicial de carga; ¿por qué se apaga después?
\end{enumerate}

\medskip
\noindent\textbf{Parte III --- El milisegundo de gloria.}
\begin{enumerate}[resume]
  \item Escribe la ley de descarga $u(t)$ a través del tubo y calcula la nueva
        \omterm{def:g12:rc-rl-circuits:tau}{constante de tiempo} $\tau'$.
  \item Calcula la corriente máxima por el tubo.
  \item Calcula la \omterm{def:g11:work-of-force:power}{potencia} máxima; compárala con la de un hervidor de \qty{2.2}{kW}.
  \item Tomando el destello por terminado en $3\tau'$, da su duración y su \omterm{def:g11:work-of-force:power}{potencia}
        media.
  \item Forma el cociente entre el tiempo de carga y el tiempo de destello: el mismo
        \omterm{def:g12:rc-rl-circuits:capacitor}{condensador}, la misma carga --- ¿qué ha cambiado, y en qué convierte eso a un
        \omterm{def:g12:rc-rl-circuits:capacitor}{condensador}: en un depósito de \omterm{def:g10:energy-conservation:energy}{energía} o en una palanca de \omterm{def:g11:work-of-force:power}{potencia}?
\end{enumerate}

\medskip
\noindent\textbf{Parte IV --- El primo de la pared.}
Un desfibrilador debe entregar $E_C = \qty{200}{J}$ a $U_0 = \qty{1500}{V}$; el tórax
y las palas presentan unos \qty{50}{\ohm}.
\begin{enumerate}[resume]
  \item Calcula la \omterm{def:g12:rc-rl-circuits:capacitor}{capacidad} necesaria.
  \item Calcula la carga almacenada.
  \item Calcula la \omterm{def:g12:rc-rl-circuits:tau}{constante de tiempo} del choque y su duración ($3\tau$); compárala
        con la del flash.
  \item Calcula la corriente máxima y la \omterm{def:g11:work-of-force:power}{potencia} máxima en el paciente.
  \item El circuito de recarga debe dejar el aparato listo (al $95\,\%$) en
        \qty{6.0}{s}: halla la \omterm{def:g11:circuits-and-power:resistance}{resistencia} de carga necesaria y la corriente inicial
        de carga --- las cifras que encargaría un proyectista.
\end{enumerate}
\end{problem}
